\cleardoublepage
\chapter{PENUTUP}
\label{chap:penutup}

\section{Rangkuman Pelaksanaan Pekerjaan}
Kerja Praktik berfokus pada tiga pekerjaan. Pekerjaan pertama adalah pengembangan query SQL dan alur ETL untuk kebutuhan campaign. Pekerjaan kedua adalah analisis deskriptif campaign financing tahun 2026. Pekerjaan ketiga adalah spatial clustering pada data log aplikasi digital.

Pekerjaan tidak berjalan sebagai tiga aktivitas yang terpisah sepenuhnya. Pemahaman dari pekerjaan query membantu penulis membaca struktur data pada analisis. Proses validasi pada campaign juga terbawa ke pekerjaan clustering. Penulis menjadi lebih terbiasa memeriksa definisi, jumlah data, duplikasi, batas data, dan kemungkinan salah interpretasi sebelum membuat output.

\section{Rangkuman Hasil Pekerjaan}
Pada pekerjaan SQL dan ETL, penulis merapikan query kompleks ke struktur CTE, memeriksa hasil join dan filter, melakukan validasi bertahap, lalu menguji integrasinya pada Talend Data Integration. Penulis juga mempelajari penjadwalan proses berkala melalui Talend Administration Center. Sebagai eksplorasi tambahan, penulis mencoba penggunaan agent pada Microsoft Copilot untuk membantu cleansing dan persiapan dataset dalam pencarian pola whitelist sehingga proses pemeriksaan tidak hanya bergantung pada eye-ball checking manual.

Pada analisis campaign financing, penulis menyusun ringkasan performa berdasarkan beberapa metrik internal pada tingkat agregat. Hasil digunakan untuk membedakan campaign yang mempunyai karakter berbeda dan menyusun bahan tindak lanjut. Penulis juga mencatat batas interpretasi karena data tidak selalu mempunyai kelompok pembanding dan detail waktu yang lengkap.

Pada spatial clustering, penulis membangun pipeline dari data log lokasi menuju stay point, clustering, reverse geocoding centroid, enrichment agregat, dan visualisasi. Hasil tersebut dapat digunakan untuk membaca kandidat area aktivitas. Pengembangan SNA belum diselesaikan sehingga ditempatkan sebagai pekerjaan lanjutan.

\section{Evaluasi Diri}
Perubahan cara berpikir terbesar selama Kerja Praktik adalah bagaimana memulai pekerjaan data dari goal, bukan dari query atau tools. Pada awalnya, penulis cenderung lebih cepat masuk ke tabel dan mencoba menghasilkan output. Setelah mengerjakan beberapa kebutuhan, penulis melihat bahwa langkah pertama justru harus menjawab apa yang ingin dicapai, keputusan apa yang akan didukung, dan bentuk output apa yang dibutuhkan.

Dari goal tersebut, penulis perlu menurunkan kebutuhan data secara bertahap. Penulis menentukan informasi atau metrik yang diperlukan, entitas yang terlibat, kolom yang dibutuhkan, tabel sumber, relasi antartabel, periode data, filter, flag, dan aturan bisnis yang digunakan. Proses ini juga mencakup pemeriksaan apakah grain antarsumber sesuai, apakah satu identifier unik atau berulang, dan apakah definisi suatu atribut benar-benar sama dengan kebutuhan bisnis.

Penulis memahami bahwa kemampuan teknis menulis SQL, menjalankan Talend, atau menggunakan agent belum cukup jika pemetaan kebutuhan datanya salah. Query dapat berhasil dijalankan tetapi tetap menghasilkan jawaban yang salah ketika kolom yang dipilih tidak mewakili konsep bisnis, \textit{join} menggandakan data, atau filter memakai definisi yang tidak sesuai. Temuan tersebut membuat penulis perlu lebih disiplin menulis goal, data requirement, sumber, asumsi, dan checkpoint validasi sebelum implementasi.

Penulis juga belajar untuk membatasi scope. Pada proyek spatial, ada banyak kemungkinan pengembangan seperti rekomendasi ekosistem, community detection, sampai influence score. Tidak semua bagian dapat diselesaikan dalam waktu Kerja Praktik. Penulis memilih menyelesaikan bagian yang dapat divalidasi dan menjelaskan bagian yang belum selesai daripada memaksakan klaim hasil.

Hal yang masih perlu diperbaiki adalah dokumentasi keputusan sejak awal. Beberapa klarifikasi dan alasan perubahan baru ditulis setelah eksperimen berjalan. Ke depan, penulis perlu membuat alur kerja yang lebih eksplisit dari goal, kebutuhan data, pemetaan kolom dan tabel, implementasi, validasi, sampai interpretasi hasil agar proses review lebih mudah.

\section{Saran}
Berdasarkan pengalaman penulis, pekerjaan data campaign akan lebih mudah ditelusuri jika definisi metrik dan aturan bisnis disimpan bersama dengan kebutuhan data. Hal ini dapat mengurangi perbedaan interpretasi antara pihak yang meminta data dan pihak yang mengolah data.

Untuk pekerjaan analitik, penulis menyarankan agar eksperimen selalu mempunyai batas interpretasi yang ditulis sejak awal. Contohnya adalah perbedaan antara analisis deskriptif dan analisis kausal, serta perbedaan antara co-presence dan hubungan sosial pada data lokasi. Dokumentasi batas tersebut membantu mencegah insight dibaca lebih jauh dari data yang tersedia.

\section{Ucapan Terima Kasih}
Penulis mengucapkan terima kasih kepada Bapak M. Rasyid Aqmar selaku Pembimbing Kerja Praktik di perusahaan atas arahan selama pelaksanaan pekerjaan. Penulis juga mengucapkan terima kasih kepada Bapak Fitra Arifiansyah, S.Kom., M.T. selaku Dosen Kerja Praktik, Mas Irsyad dan Mbak Sekar, serta rekan-rekan di lingkungan kerja dari tim ataupun departemen lain yang membantu penulis selama Kerja Praktik.

Bantuan dan pembelajaran yang diberikan mencakup pemahaman proses bisnis, data, penggunaan tools, komunikasi, kolaborasi, adaptasi, serta cara bekerja di lingkungan profesional. Pengalaman tersebut membantu penulis memahami bahwa pekerjaan data tidak hanya membutuhkan kemampuan teknis, tetapi juga kemampuan untuk memahami kebutuhan, menyampaikan hasil, dan bekerja bersama pihak dengan tanggung jawab yang berbeda.
